\section{Killing the two output poles does not give an invariant raw domain}
\label{sec:pole-free-noninvariance}

The noninvariance in MA-MO-007 has a precise further refinement.
One may require the first raw output to have both polar coefficients
zero. That restriction is a closed subspace of finite codimension in
the additive Schwartz source, but it is not preserved by the raw
action. We prove this with genuine compact additive inputs and two
successive original source compositions. No estimate for the norm of
the once-defined restricted operator is inferred from this result.

Keep the rational source place set \(\{\infty\}\), the real even
cutoff \(v(t)=\phi(\log|t|)\), \(v(0)=1\), and the actual maps
of MA-MO-007. Thus \(v\in C_c^\infty(\mathbb R)\) is one near
zero. Write \(\mathscr S_{\rm ev}\) for the even Schwartz functions.
On that space \(\mathfrak F^*=\mathfrak F\) and
\(\mathfrak F^2=1\). Define
\[
 \begin{aligned}
 L&=1-M_v-\mathfrak F M_v\mathfrak F,\\
 \ell_0(f)&=\int_{\mathbb R}v(t)f(t)\,dt,\\
 \ell_1(f)&=\int_{\mathbb R}v(t)\mathfrak Ff(t)\,dt
             =\int_{\mathbb R}\mathfrak Fv(t)f(t)\,dt,\\
 \mathscr U_0&=\{f\in\mathscr S_{\rm ev}:f(0)=\int f=0\},\\
 \mathscr U_{\rm pf}&=\{f\in\mathscr U_0:
                                      \ell_0(f)=\ell_1(f)=0\}.
 \end{aligned}
 \tag{MA-MO.89}
\]
The equality of the two integrals for \(\ell_1\) is ordinary
Fourier Fubini on Schwartz functions. All functionals in this display
are continuous in the Schwartz topology. Let
\(D=P_+-P_-\) denote the actual source matrix.

\paragraph{MA-MO-008 (exact pole-free source and failure at the second iterate).}
The subspace \(\mathscr U_{\rm pf}\) is closed of codimension four
in \(\mathscr S_{\rm ev}\), and of codimension two in
\(\mathscr U_0\). For \(f\in\mathscr U_0\), the actual raw output
\(D i_+f\) belongs to the inherited narrow Hilbert product if and
only if \(f\in\mathscr U_{\rm pf}\). In that case it is a diagonal
pair in \(\mathcal E\oplus\mathcal E\).

Nevertheless there is a real even
\(f_0\in C_c^\infty(\mathbb R\setminus\{0\})\cap
\mathscr U_{\rm pf}\) such that, with \(f_1=Lf_0\) and
\(f_2=L^2f_0\),
\[
 \begin{gathered}
 f_1\in\mathscr U_0,\qquad \ell_0(f_1)=1,
                          \qquad \int f_2=-1,\\
 D i_+f_0=i_+f_1=(\Theta f_1,\Theta f_1)
                                     \in\mathcal E\oplus\mathcal E,\\
 D^2 i_+f_0=i_+f_2,\qquad
                         \Theta f_2\notin\mathcal H_a
                                     \quad(0<a<1/2).
 \end{gathered}
 \tag{MA-MO.90}
\]
Both iterates in this formula are defined in the original smooth
source product. Thus restricting once to pole-free outputs does not
produce an invariant raw domain.

\begin{proof}
\emph{The exact raw action and its polar functionals.}
The operator \(L\) preserves even Schwartz space. As an operator on
the ordinary additive even \(L^2\) space, it is bounded and
self-adjoint, since \(v\) is real and bounded and the Fourier
transform is unitary and involutive there. It commutes with
\(\mathfrak F\). Direct Fourier evaluation gives, for every
\(f\in\mathscr S_{\rm ev}\),
\[
                 (Lf)(0)=-\ell_1(f),\qquad
                        \int Lf=-\ell_0(f).
 \tag{MA-MO.91}
\]
Indeed \(v(0)=1\) cancels the first two values in \(Lf(0)\).
For the integral, \(\int\mathfrak F(v\mathfrak Ff)
=(v\mathfrak Ff)(0)=\int f\), leaving \(-\ell_0(f)\).

Formula (MA-MO.85) is exactly
\(p_+i_+=1+L\). For \(f\in\mathscr U_0\), the source vector
\(i_+f\) is diagonal and \(P_-i_+f=i_+f\). Therefore
\[
                       D i_+f=i_+Lf
                              \qquad(f\in\mathscr U_0).
 \tag{MA-MO.92}
\]
This identity uses the actual source maps, with no idempotence
assumption. By MA-MO-007, a source common-sum vector belongs to
the narrow Hilbert product exactly when its polar difference
vanishes. The two additive polar moments of \(Lf\) are
\(-\ell_1(f)\) and \(-\ell_0(f)\). By (MA-MO.84), the
difference of the first and second components of \(i_+Lf\)
is therefore \(\ell_1(f)e^{x/2}-\ell_0(f)e^{-x/2}\).
Thus the claimed condition is necessary.
It is sufficient because both polar moments of \(Lf\) then
vanish, and full Poisson summation makes
\(\Theta Lf\in\mathcal E\). Here \(Lf\) is an additive
Schwartz function; it need not be flat at zero, and its
centered preimage is not asserted to belong to \(\mathcal E\).

\emph{The stated finite codimensions.}
The three smooth functions
\(1,v,\mathfrak Fv\) are linearly independent on
\(\mathbb R\setminus\{0\}\). To see this, an identity
\(b_0+b_1v+b_2\mathfrak Fv=0\) outside \(\operatorname{supp}v\)
first gives \(b_0=0\) by Schwartz decay. If \(b_2\ne0\),
the entire function \(\mathfrak Fv\) vanishes on an open
interval, hence vanishes identically, contradicting \(v=1\)
near zero. Thus \(b_2=0\), and then \(b_1=0\).
The transform is entire because \(v\) has compact support;
differentiating its integral proves this directly.

These three functions represent the integral, \(\ell_0\),
and \(\ell_1\). Their independence as functionals on even
compact tests away from zero follows by testing with even
bumps, since all three are even and continuous. Evaluation
at zero is independent of them on \(\mathscr S_{\rm ev}\):
a distributional relation would express a multiple of
\(\delta_0\) as a smooth function; testing away from zero
and using the independence just proved forces the smooth
part to vanish, and then the multiple of \(\delta_0\) is
zero too. Thus the four continuous functionals are independent.
Their joint map has rank four, and its kernel is
\(\mathscr U_{\rm pf}\). The first two have rank two,
so the asserted codimensions and closedness follow.

\emph{The additional functional is not a consequence of the first three.}
Put
\[
                      W=Lv=v-v^2-\mathfrak F(v\mathfrak Fv).
 \tag{MA-MO.93}
\]
This is a real even Schwartz function. We claim
\[
                  W\notin\operatorname{span}\{1,v,\mathfrak Fv\}
                  \quad\hbox{on }\mathbb R\setminus\{0\}.
 \tag{MA-MO.94}
\]
Suppose otherwise. Outside \(\operatorname{supp}v\), the
purported identity is
\(-\mathfrak F(v\mathfrak Fv)=b_0+b_2\mathfrak Fv\).
Both Fourier transforms here are Schwartz, since their
preimages are compact smooth functions. Sending \(t\) to
positive infinity gives \(b_0=0\). Both transforms are also
entire, so equality on an exterior interval implies equality
everywhere. Fourier inversion in the even sector yields
\[
                         v\mathfrak Fv=-b_2v.
 \tag{MA-MO.95}
\]
On a neighborhood of zero where \(v=1\), this says
\(\mathfrak Fv=-b_2\). The entire function \(\mathfrak Fv\)
would therefore be constant everywhere, and its Schwartz
decay would force it to be zero. Fourier injectivity then
forces \(v=0\), a contradiction. This proves (MA-MO.94).

\emph{An actual finite-bump construction.}
For clarity, the required annihilating test can be obtained
from a finite linear system. The four real even functions
\[
                   k_1=1,\quad k_2=v,\quad
                   k_3=\mathfrak Fv,\quad k_4=W
 \tag{MA-MO.96}
\]
are linearly independent on \((0,\infty)\), by the preceding
arguments and evenness. Hence there are distinct positive
points \(t_1,\ldots,t_4\) for which the evaluation matrix
\((k_i(t_j))_{i,j=1}^4\) is invertible. Otherwise all evaluation
vectors would span a proper subspace of \(\mathbb R^4\),
and a nonzero linear functional vanishing on their span
would give a forbidden relation among the four functions.

Choose an even nonnegative \(\eta\in C_c^\infty((-1,1))\)
of integral one, and define, for sufficiently small \(\epsilon>0\),
\[
 \begin{aligned}
 \beta_j^\epsilon(t)
   &=\frac{1}{2\epsilon}
       \left[\eta\left(\frac{t-t_j}{\epsilon}\right)
              +\eta\left(\frac{t+t_j}{\epsilon}\right)\right],\\
 M_{ij}(\epsilon)&=\int_{\mathbb R}k_i(t)\beta_j^\epsilon(t)\,dt.
 \end{aligned}
 \tag{MA-MO.97}
\]
These are even compact smooth bumps of mass one with disjoint
supports away from zero. The matrix tends to the invertible
evaluation matrix, so is invertible for all sufficiently
small \(\epsilon\). Fix one such value. Solve
\(M(\epsilon)c=(0,0,0,1)^{\mathsf T}\), and set
\(f_0=\sum_{j=1}^4c_j\beta_j^\epsilon\). The coefficients
and the function are real. By construction
\[
 f_0(0)=\int f_0=\ell_0(f_0)=\ell_1(f_0)=0,
                         \qquad \int Wf_0=1.
 \tag{MA-MO.98}
\]
This is an actual compact additive test, not a distribution
or an assumed vector of a completed quotient.

\emph{Two original compositions and the exit at the second one.}
By (MA-MO.91), \(f_1=Lf_0\) is even Schwartz with
\(f_1(0)=\int f_1=0\). Since the data are real, the
self-adjointness of \(L\), or direct Fourier Fubini, gives
\[
 \ell_0(f_1)=\int vLf_0=\int (Lv)f_0=1,
               \qquad \int f_2=-\ell_0(f_1)=-1.
 \tag{MA-MO.99}
\]
The first source action is (MA-MO.92), and its output is the
genuine diagonal \(i_+f_1\). Applying (MA-MO.92) again is
legitimate because \(f_1\in\mathscr U_0\). It gives the
second identity in (MA-MO.90). The first component
\(\Theta f_2\) now has left polar term \(-e^{-x/2}\),
by the full Poisson formula and \(\int f_2=-1\). Its other
polar term is a multiple of \(e^{x/2}\); the remainder is
rapid at the left. Hence \(\Theta f_2\notin\mathcal H_a\)
for \(0<a<1/2\). This proves every assertion.
\end{proof}

\paragraph{What the result does and does not repair.}
The pole-free source conditions are exact and finite in number,
but applying them once is insufficient for a raw invariant
realization: the explicitly constructed first output is a valid
diagonal smooth vector whose next original output leaves the
Hilbert target. A larger family of conditions, or a chosen
polar correction after each action, defines a different domain
or operator and needs a separate proof.

This is not a sequence with unbounded one-step Hilbert norm,
and it is not a nonclosability proof for the restriction from
\(\mathscr U_{\rm pf}\). In particular boundedness of \(L\)
on ordinary additive \(L^2\) does not establish boundedness
of its theta-conjugated action for the inherited
\(H_{\Sigma,a}\) norm. No density or closure description
for the pole-free theta-image domain in that norm is asserted.
The PR quotient returns and the RS continued operators retain
their distinct, already proved scopes.

The source identities used here are exactly the approximate-section
and Poisson maps of MA-MO.84--.85, with Meyer (2005),
\texttt{lem:Hilm\_p\_embed}, the \(p_+^S i_+\) display,
the defining matrices, and \texttt{eq:Poisson\_summation},
author-manuscript pp.~39--41. The inherited finite polar
representation and finite-length spectral \(K_0\) data do
not make an arbitrary kernel of polar functionals invariant.
The cote~1 nuclear-trace placement, scans~079--080 and
111--112 through ALPOGE\_BRIDGES BR-TR-01/02, still requires
actual maps on the claimed domains before iteration or trace
identities. No new primary-source audit is claimed beyond
the precise MO004--007 checks, and no zero-location premise
enters the present argument.
